%-----------------------------------------------------------------------------------------------------------------------------------------------%
%	The MIT License (MIT)
%
%	Copyright (c) 2015 Jan Küster
%
%	Permission is hereby granted, free of charge, to any person obtaining a copy
%	of this software and associated documentation files (the "Software"), to deal
%	in the Software without restriction, including without limitation the rights
%	to use, copy, modify, merge, publish, distribute, sublicense, and/or sell
%	copies of the Software, and to permit persons to whom the Software is
%	furnished to do so, subject to the following conditions:
%	
%	THE SOFTWARE IS PROVIDED "AS IS", WITHOUT WARRANTY OF ANY KIND, EXPRESS OR
%	IMPLIED, INCLUDING BUT NOT LIMITED TO THE WARRANTIES OF MERCHANTABILITY,
%	FITNESS FOR A PARTICULAR PURPOSE AND NONINFRINGEMENT. IN NO EVENT SHALL THE
%	AUTHORS OR COPYRIGHT HOLDERS BE LIABLE FOR ANY CLAIM, DAMAGES OR OTHER
%	LIABILITY, WHETHER IN AN ACTION OF CONTRACT, TORT OR OTHERWISE, ARISING FROM,
%	OUT OF OR IN CONNECTION WITH THE SOFTWARE OR THE USE OR OTHER DEALINGS IN
%	THE SOFTWARE.
%	
%
%-----------------------------------------------------------------------------------------------------------------------------------------------%


%============================================================================%
%
%	DOCUMENT DEFINITION
%
%============================================================================%

%we use article class because we want to fully customize the page and don't use a cv template
\documentclass[10pt,A4]{article}	


%----------------------------------------------------------------------------------------
%	ENCODING
%----------------------------------------------------------------------------------------

%we use utf8 since we want to build from any machine
\usepackage[utf8]{inputenc}		

%----------------------------------------------------------------------------------------
%	LOGIC
%----------------------------------------------------------------------------------------

% provides \isempty test
\usepackage{xifthen}

%----------------------------------------------------------------------------------------
%	FONT
%----------------------------------------------------------------------------------------

% some tex-live fonts - choose your own

%\usepackage[defaultsans]{droidsans}
%\usepackage[default]{comfortaa}
%\usepackage{cmbright}
%\usepackage[default]{raleway}
%\usepackage{fetamont}
%\usepackage[default]{gillius}
%\usepackage[light,math]{iwona}
%\usepackage[thin]{roboto} 

% set font default
\renewcommand*\familydefault{\sfdefault} 	
\usepackage[T1]{fontenc}

% more font size definitions
\usepackage{moresize}		


%----------------------------------------------------------------------------------------
%	PAGE LAYOUT  DEFINITIONS
%----------------------------------------------------------------------------------------

%debug page outer frames
%\usepackage{showframe}			


%define page styles using geometry
\usepackage[a4paper]{geometry}		

% for example, change the margins to 2 inches all round
\geometry{top=1.25cm, bottom=-.6cm, left=1.5cm, right=1.5cm} 	

%use customized header
\usepackage{fancyhdr}				
\pagestyle{fancy}

%less space between header and content
\setlength{\headheight}{-5pt}		


%customize entries left, center and right
%\lhead{}
%\chead{}
%\rhead{}


%indentation is zero
\setlength{\parindent}{0mm}

%----------------------------------------------------------------------------------------
%	TABLE /ARRAY DEFINITIONS
%---------------------------------------------------------------------------------------- 

%for layouting tables
\usepackage{multicol}			
\usepackage{multirow}

%extended aligning of tabular cells
\usepackage{array}

\newcolumntype{x}[1]{%
>{\raggedleft\hspace{0pt}}p{#1}}%


%----------------------------------------------------------------------------------------
%	GRAPHICS DEFINITIONS
%---------------------------------------------------------------------------------------- 

%for header image
\usepackage{graphicx}

%for floating figures
\usepackage{wrapfig}
\usepackage{float}
%\floatstyle{boxed} 
%\restylefloat{figure}

%for drawing graphics		
\usepackage{tikz}				
\usetikzlibrary{shapes, backgrounds,mindmap, trees}


%----------------------------------------------------------------------------------------
%	Color DEFINITIONS
%---------------------------------------------------------------------------------------- 

\usepackage{color}

%accent color
\definecolor{sectcol}{RGB}{0,0,180}

%dark background color
\definecolor{bgcol}{RGB}{110,110,110}

%light background / accent color
\definecolor{softcol}{RGB}{225,225,225}


%----------------------------------------------------------------------------------------
%	Other Packages
%---------------------------------------------------------------------------------------- 

% for \href
\definecolor{linkcolor}{RGB}{0,50,150}
\usepackage[colorlinks=true,urlcolor=linkcolor,linkcolor=linkcolor]{hyperref}

%============================================================================%
%
%
%	DEFINITIONS
%
%
%============================================================================%

%----------------------------------------------------------------------------------------
% 	HEADER
%----------------------------------------------------------------------------------------

% remove top header line
\renewcommand{\headrulewidth}{0pt} 

%remove bottom header line
\renewcommand{\footrulewidth}{0pt}	  	

%remove pagenum
\renewcommand{\thepage}{}	

%remove section num		
\renewcommand{\thesection}{}			

%----------------------------------------------------------------------------------------
% 	ARROW GRAPHICS in Tikz
%----------------------------------------------------------------------------------------

% a six pointed arrow pointing to the left
\newcommand{\tzlarrow}{(0,0) -- (0.2,0) -- (0.3,0.2) -- (0.2,0.4) -- (0,0.4) -- (0.1,0.2) -- cycle;}	

% include the left arrow into a tikz picture
% param1: fill color
%
\newcommand{\larrow}[1]
{\begin{tikzpicture}[scale=0.58]
	 \filldraw[fill=#1!100,draw=#1!100!black]  \tzlarrow
 \end{tikzpicture}
}

% a six pointed arrow pointing to the right
\newcommand{\tzrarrow}{ (0,0.2) -- (0.1,0) -- (0.3,0) -- (0.2,0.2) -- (0.3,0.4) -- (0.1,0.4) -- cycle;}

% include the right arrow into a tikz picture
% param1: fill color
%
\newcommand{\rarrow}
{\begin{tikzpicture}[scale=0.7]
	\filldraw[fill=sectcol!100,draw=sectcol!100!black] \tzrarrow
 \end{tikzpicture}
}



%----------------------------------------------------------------------------------------
%	custom sections
%----------------------------------------------------------------------------------------

% create a coloured box with arrow and title as cv section headline
% param 1: section title
%
\newcommand{\cvsection}[1]
{
	\begin{center}
		\large\textcolor{sectcol}{\textbf{#1}}
	\end{center}
}

%create a coloured arrow with title as cv meta section section
% param 1: meta section title
%
\newcommand{\metasection}[2]
{
%\begin{tabular*}{1\textwidth}{r r}
\footnotesize{#2} \hspace*{\fill} \footnotesize{#1}\\[1pt]
%\end{tabular*}
}

%----------------------------------------------------------------------------------------
%	 CV EVENT
%----------------------------------------------------------------------------------------

% creates a stretched box as cv entry headline followed by some paragraphs about 
% the work you did
% param 1:	event time i.e. 2014 or 2011-2014 etc.
% param 2:	event name (what did you do?)
% param 3:	institution (where did you work / study)
% param 4:	list of paragraphs
%
\newcommand{\cvevent}[4]
{

\begin{tabular*}{1\textwidth}{p{13.6cm}  x{3.9cm}}
	\textbf{#2} - \textcolor{bgcol}{#3} &   \vspace{2.5pt}\textcolor{sectcol}{#1}
\end{tabular*}

\vspace{-8pt}
\textcolor{softcol}{\hrule}
\vspace{6pt}

	\foreach \desc in {#4}{
		$\cdot$ \desc\\[3pt]
	}
	
\vspace{3pt}
}

% creates a stretched box as 
\newcommand{\cveventmeta}[2]
{
	\mbox{\mystrut \hspace{87pt}\textit{#1}}\\
	#2
}

%----------------------------------------------------------------------------------------
% CUSTOM STRUT FOR EMPTY BOXES
%----------------------------------------- -----------------------------------------------
\newcommand{\mystrut}{\rule[-.3\baselineskip]{0pt}{\baselineskip}}

%----------------------------------------------------------------------------------------
% CUSTOM LOREM IPSUM
%----------------------------------------------------------------------------------------
\newcommand{\lorem}
{Lorem ipsum dolor sit amet, consectetur adipiscing elit. Donec a diam lectus.}



%============================================================================%
%
%
%
%	DOCUMENT CONTENT
%
%
%
%============================================================================%
\begin{document}

\pagestyle{fancy}

\vspace{-8pt}
\begin{center}
    {\HUGE \textsc{Boudrouss Réda} \textcolor{sectcol}{\rule[-1mm]{1mm}{0.9cm}} \textsc{CV}}\\[2pt]
    {\small Élève ingénieur en informatique -- Développement logiciel}
\end{center}

\vspace{6pt}

%----------------------------------------
%   MÉTADONNÉES
%----------------------------------------
\metasection{\href{https://github.com/rboudrouss}{github.com/rboudrouss}}{}
\metasection{\href{https://www.linkedin.com/in/rboudrouss}{linkedin.com/in/rboudrouss}}{}
\metasection{\href{mailto:contact@rboud.com}{contact@rboud.com}}{}
\metasection{\href{https://rboud.com}{rboud.com}}{}
\metasection{+33 (0)7 55 78 85 16}{\textbf{Langues :} Français (natif), Anglais (professionnel), Arabe (conversationnel)}

\vspace{-15pt}
\textcolor{softcol}{\hrule}

\normalsize

%============================================================================%
% RÉSUMÉ
%============================================================================%

\cvsection{Résumé}

Étudiant en Master \textbf{Sciences et Technologies du Logiciel} à Sorbonne Université, passionné par le \textbf{développement logiciel}, la conception de systèmes robustes et les technologies modernes. J'apprécie autant le développement de nouvelles fonctionnalités que l'amélioration de la qualité, des performances et de la maintenabilité d'applications existantes. Curieux et polyvalent, je m'intéresse aussi bien au développement back-end qu'aux outils, aux infrastructures et aux problématiques de génie logiciel.

%============================================================================%
% FORMATION
%============================================================================%

\cvsection{Formation}

\cvevent{2024 -- 2026}{Master Sciences et Technologies du Logiciel}{Sorbonne Université}{
    {Approfondissement en génie logiciel, compilation, langages de programmation, systèmes, concurrence et architectures distribuées.},
    {Conception et implémentation de compilateurs, d'interpréteurs et d'outils logiciels, avec une attention portée à la qualité, aux performances et à la robustesse.},
    {Nombreux projets en OCaml, C/C++, Java, Python et Rust autour du développement logiciel, des systèmes et des environnements d'exécution.}
}

\cvevent{2020 -- 2024}{Licence Informatique}{Sorbonne Université}{
    {Formation en algorithmique, architecture des ordinateurs, systèmes d'exploitation, réseaux, bases de données, programmation orientée objet et fonctionnelle.},
    {Réalisation de nombreux projets individuels et collectifs dans différents langages et environnements de développement.}
}

%============================================================================%
% EXPÉRIENCES
%============================================================================%

\vspace{-20pt}
\cvsection{Expériences}

\cvevent{2025 -- 2026}{Quality Analyst Engineer}{Hexaglobe}{
    {Conception et développement d'une stratégie complète de tests (unitaires, intégration et end-to-end) sur plusieurs applications, représentant plus de \textbf{3 000 cas de test}.},
    {Développement d'outils internes pour automatiser la production de tests et améliorer la productivité des équipes grâce à des workflows assistés par IA.},
    {Mise en place de pipelines \textbf{CI/CD GitLab}, mutualisation des composants de build et amélioration des processus de livraison logicielle.},
    {Développement d'outils de validation multimédia et amélioration de la qualité des produits de streaming vidéo.}
}

\cvevent{2023 -- 2026}{Responsable Technique}{Association Play Sorbonne Université}{
    {Conception, développement et administration de l'ensemble de l'infrastructure technique de l'association (Docker, CI/CD, serveurs, Google Workspace).},
    {Développement et maintenance du site web de l'association (\textbf{9 000 visiteurs uniques} en 2025) avec optimisation des performances et du référencement.},
    {Conception d'une application web mobile \textbf{local-first} avec authentification sécurisée, architecture back-end optimisée et base de données adaptée à de fortes charges.}
}

\cvevent{2024 -- 2025}{Président}{Association Play Sorbonne Université}{
    {Pilotage de l'association, coordination d'équipes techniques et organisationnelles, suivi des projets et définition des priorités.},
    {Encadrement de plus de \textbf{100 bénévoles} dans l'organisation d'un événement accueillant plus de \textbf{10 000 visiteurs}.}
}

%============================================================================%
% PROJETS
%============================================================================%

\vspace{-20pt}
\cvsection{Projets}

$\cdot$
\href{https://github.com/rboudrouss/mopsa-emcc}{\textbf{Mopsa-WASM}} :
Portage d'un logiciel complexe vers \textbf{WebAssembly} et développement d'une interface web. Travail sur une importante base de code en \textbf{OCaml, C et C++}.

\vspace{2pt}

$\cdot$
\href{https://github.com/rboudrouss/reactant-analyzer}{\textbf{Reactant Analyzer}} :
Développement d'un outil en \textbf{Rust} pour analyser des applications React, avec un fort accent sur l'architecture logicielle, les performances et la qualité du code.

\vspace{2pt}

$\cdot$
Développement de nombreux projets universitaires en \textbf{C, C++, Java, Python, Rust, OCaml et TypeScript}, couvrant des domaines tels que les systèmes, les compilateurs, les applications web, les architectures distribuées et les outils de développement.

%--------------------------------------------------------------------------------------------------
% FOOTER
%--------------------------------------------------------------------------------------------------


\end{document}